\documentclass[10pt,a4paper]{amsart}
\usepackage{amsmath,amssymb,amsthm,graphicx}
% Retained preamble for the Stacks Project chapter {CH} (shared project preamble).
% The chapter file pinned-original-source.tex only carries \input{{preamble}} on its first line,
% so no preamble text lives inside the chapter itself.  The Stacks Project uses one single shared
% preamble for every chapter: the snapshot stages it three times (sources/stacks-chow/preamble.tex,
% sources/stacks-etale-cohomology/preamble.tex and focused-seeds/AGP/sources/stacks-snapshot-20261011/
% preamble.tex) and those copies are identical in size.  The block below is copied mechanically,
% line for line, from sources/stacks-chow/preamble.tex lines 14-18 and 154-217: the Xy-pic setup and
% the native theorem declarations together with the mathematical macros.  Nothing was added, removed
% or redefined; the wrapper class and the amsmath/amssymb/amsthm/graphicx base packages are supplied
% by the assembler.
% For commutative diagrams we use Xy-pic
% Chapter: stacks-spaces-more-morphisms.tex -- More on Morphisms of Spaces.
\usepackage[all]{xy}

% We use 2cell for 2-commutative diagrams.
\xyoption{2cell}

% Theorem environments.
%
\theoremstyle{plain}
\newtheorem{theorem}[subsection]{Theorem}
\newtheorem{proposition}[subsection]{Proposition}
\newtheorem{lemma}[subsection]{Lemma}

\theoremstyle{definition}
\newtheorem{definition}[subsection]{Definition}
\newtheorem{example}[subsection]{Example}
\newtheorem{exercise}[subsection]{Exercise}
\newtheorem{situation}[subsection]{Situation}

\theoremstyle{remark}
\newtheorem{remark}[subsection]{Remark}
\newtheorem{remarks}[subsection]{Remarks}

\numberwithin{equation}{subsection}

% Macros
%
\def\lim{\mathop{\mathrm{lim}}\nolimits}
\def\colim{\mathop{\mathrm{colim}}\nolimits}
\def\Spec{\mathop{\mathrm{Spec}}}
\def\Hom{\mathop{\mathrm{Hom}}\nolimits}
\def\Ext{\mathop{\mathrm{Ext}}\nolimits}
\def\SheafHom{\mathop{\mathcal{H}\!\mathit{om}}\nolimits}
\def\SheafExt{\mathop{\mathcal{E}\!\mathit{xt}}\nolimits}
\def\Sch{\mathit{Sch}}
\def\Mor{\mathop{\mathrm{Mor}}\nolimits}
\def\Ob{\mathop{\mathrm{Ob}}\nolimits}
\def\Sh{\mathop{\mathit{Sh}}\nolimits}
\def\NL{\mathop{N\!L}\nolimits}
\def\CH{\mathop{\mathrm{CH}}\nolimits}
\def\proetale{{pro\text{-}\acute{e}tale}}
\def\etale{{\acute{e}tale}}
\def\QCoh{\mathit{QCoh}}
\def\Ker{\mathop{\mathrm{Ker}}}
\def\Im{\mathop{\mathrm{Im}}}
\def\Coker{\mathop{\mathrm{Coker}}}
\def\Coim{\mathop{\mathrm{Coim}}}

% Boxtimes
%
\DeclareMathSymbol{\boxtimes}{\mathbin}{AMSa}{"02}

%
% Macros for moduli stacks/spaces
%
\def\QCohstack{\mathcal{QC}\!\mathit{oh}}
\def\Cohstack{\mathcal{C}\!\mathit{oh}}
\def\Spacesstack{\mathcal{S}\!\mathit{paces}}
\def\Quotfunctor{\mathrm{Quot}}
\def\Hilbfunctor{\mathrm{Hilb}}
\def\Curvesstack{\mathcal{C}\!\mathit{urves}}
\def\Polarizedstack{\mathcal{P}\!\mathit{olarized}}
\def\Complexesstack{\mathcal{C}\!\mathit{omplexes}}
% \Pic is the operator that assigns to X its picard group, usage \Pic(X)
% \Picardstack_{X/B} denotes the Picard stack of X over B
% \Picardfunctor_{X/B} denotes the Picard functor of X over B
\def\Pic{\mathop{\mathrm{Pic}}\nolimits}
\def\Picardstack{\mathcal{P}\!\mathit{ic}}
\def\Picardfunctor{\mathrm{Pic}}
\def\Deformationcategory{\mathcal{D}\!\mathit{ef}}
\newenvironment{reference}{}{}
\newenvironment{slogan}{}{}
\newenvironment{history}{}{}
\usepackage{cleveref}
\crefname{originallemma}{original lemma}{original lemmas}
\Crefname{originallemma}{Original lemma}{Original lemmas}
\makeatletter
\newlabel{lemma-unramified-formally-unramified}{{lemma-unramified-formally-unramified}{source}{}{sourcepointer-1}{}}
\newlabel{lemma-unramified-formally-unramified@cref}{{[originallemma][0][]lemma-unramified-formally-unramified}{[1][0][]source}{}{}{}}
\makeatother
\crefname{originallemma}{original lemma}{original lemmas}
\Crefname{originallemma}{Original lemma}{Original lemmas}
\makeatletter
\newlabel{spaces-morphisms-lemma-diagonal-unramified-morphism}{{spaces-morphisms-lemma-diagonal-unramified-morphism}{source}{}{sourcepointer-2}{}}
\newlabel{spaces-morphisms-lemma-diagonal-unramified-morphism@cref}{{[originallemma][0][]spaces-morphisms-lemma-diagonal-unramified-morphism}{[1][0][]source}{}{}{}}
\makeatother
\crefname{originallemma}{original lemma}{original lemmas}
\Crefname{originallemma}{Original lemma}{Original lemmas}
\makeatletter
\newlabel{spaces-morphisms-lemma-locally-quasi-finite-separated-representable}{{spaces-morphisms-lemma-locally-quasi-finite-separated-representable}{source}{}{sourcepointer-3}{}}
\newlabel{spaces-morphisms-lemma-locally-quasi-finite-separated-representable@cref}{{[originallemma][0][]spaces-morphisms-lemma-locally-quasi-finite-separated-representable}{[1][0][]source}{}{}{}}
\makeatother
\crefname{originallemma}{original lemma}{original lemmas}
\Crefname{originallemma}{Original lemma}{Original lemmas}
\makeatletter
\newlabel{spaces-morphisms-lemma-monomorphism-loc-finite-type-loc-quasi-finite}{{spaces-morphisms-lemma-monomorphism-loc-finite-type-loc-quasi-finite}{source}{}{sourcepointer-4}{}}
\newlabel{spaces-morphisms-lemma-monomorphism-loc-finite-type-loc-quasi-finite@cref}{{[originallemma][0][]spaces-morphisms-lemma-monomorphism-loc-finite-type-loc-quasi-finite}{[1][0][]source}{}{}{}}
\makeatother
\crefname{originallemma}{original lemma}{original lemmas}
\Crefname{originallemma}{Original lemma}{Original lemmas}
\makeatletter
\newlabel{spaces-morphisms-lemma-universally-injective}{{spaces-morphisms-lemma-universally-injective}{source}{}{sourcepointer-5}{}}
\newlabel{spaces-morphisms-lemma-universally-injective@cref}{{[originallemma][0][]spaces-morphisms-lemma-universally-injective}{[1][0][]source}{}{}{}}
\makeatother
\title{Four equivalent conditions for an unramified monomorphism of algebraic spaces}
\author{The Stacks Project Authors}
\date{stacks-project:spaces-more-morphisms.tex:local-snapshot-20261011; GNU FDL 1.2 or later}
\begin{document}
\maketitle
\noindent\textit{Editorial note.} Statement span is lines 2866-2880 of sources/stacks-spaces-more-morphisms/pinned-original-source.tex; proof span is lines 2882-2903, the complete author proof (20 body lines).\par\medskip
\noindent\textit{Editorial note.} Navigation stubs: 5 cited-but-not-included labels are supplied as non-load-bearing pointers with their pinned tags where the snapshot tag index has them.\par\medskip
\noindent\textit{Editorial note.} extra\_packages is empty on purpose: the retained preamble already loads xy with options and a second declaration clashes in the assembler wrapper.\par\medskip
\noindent\textit{Editorial note.} Text is retained verbatim, including the author's own typos and placeholders.\par\medskip
\noindent\textit{Allowed standard prerequisites (editorial).}\begin{itemize}
\item Explicit prerequisite boundary. Unramified morphisms of algebraic spaces, monomorphisms, universal injectivity, and formal unramifiedness; the comparisons are retained.
\end{itemize}\medskip
\ifdefined\hypertarget\hypertarget{sourcepointer-1}{}\fi
\noindent\textit{Source reference.} Original lemma lemma-unramified-formally-unramified: sources/stacks-spaces-more-morphisms/pinned-original-source.tex, source lines 2829--2837. Cited by the retained statement or proof but not part of the collected spans; supplied only as a navigation pointer. The argument does not depend on its content.\par\medskip
\ifdefined\hypertarget\hypertarget{sourcepointer-2}{}\fi
\noindent\textit{Source reference.} Original lemma spaces-morphisms-lemma-diagonal-unramified-morphism: sources/stacks-spaces-more-morphisms/pinned-original-source.tex, source lines 2903--2903. Cited by the retained statement or proof but not part of the collected spans; supplied only as a navigation pointer. The argument does not depend on its content.\par\medskip
\ifdefined\hypertarget\hypertarget{sourcepointer-3}{}\fi
\noindent\textit{Source reference.} Original lemma spaces-morphisms-lemma-locally-quasi-finite-separated-representable: sources/stacks-spaces-more-morphisms/pinned-original-source.tex, source lines 2903--2903. Cited by the retained statement or proof but not part of the collected spans; supplied only as a navigation pointer. The argument does not depend on its content.\par\medskip
\ifdefined\hypertarget\hypertarget{sourcepointer-4}{}\fi
\noindent\textit{Source reference.} Original lemma spaces-morphisms-lemma-monomorphism-loc-finite-type-loc-quasi-finite: sources/stacks-spaces-more-morphisms/pinned-original-source.tex, source lines 2903--2903. Cited by the retained statement or proof but not part of the collected spans; supplied only as a navigation pointer. The argument does not depend on its content.\par\medskip
\ifdefined\hypertarget\hypertarget{sourcepointer-5}{}\fi
\noindent\textit{Source reference.} Original lemma spaces-morphisms-lemma-universally-injective: sources/stacks-spaces-more-morphisms/pinned-original-source.tex, source lines 2903--2903. Cited by the retained statement or proof but not part of the collected spans; supplied only as a navigation pointer. The argument does not depend on its content.\par\medskip
\begin{lemma}
\label{lemma-universally-injective-unramified}
Let $S$ be a scheme.
Let $f : X \to Y$ be a morphism of algebraic spaces over $S$.
The following are equivalent:
\begin{enumerate}
\item $f$ is unramified and a monomorphism,
\item $f$ is unramified and universally injective,
\item $f$ is locally of finite type and a monomorphism,
\item $f$ is universally injective, locally of finite type, and
formally unramified.
\end{enumerate}
Moreover, in this case $f$ is also representable, separated, and
locally quasi-finite.
\end{lemma}

\begin{proof}
We have seen in
Lemma \ref{lemma-unramified-formally-unramified}
that being formally unramified and locally of finite type is the same thing
as being unramified.
Hence (4) is equivalent to (2).
A monomorphism is certainly formally unramified hence (3) implies (4).
It is clear that (1) implies (3). Finally, if (2) holds, then
$\Delta : X \to X \times_Y X$ is both an open immersion
(Morphisms of Spaces, Lemma
\ref{spaces-morphisms-lemma-diagonal-unramified-morphism})
and surjective
(Morphisms of Spaces, Lemma
\ref{spaces-morphisms-lemma-universally-injective})
hence an isomorphism, i.e., $f$ is a monomorphism. In this way we see that
(2) implies (1).
Finally, we see that $f$ is representable, separated, and locally
quasi-finite by
Morphisms of Spaces, Lemmas
\ref{spaces-morphisms-lemma-monomorphism-loc-finite-type-loc-quasi-finite} and
\ref{spaces-morphisms-lemma-locally-quasi-finite-separated-representable}.
\end{proof}

\end{document}
